% A vertical command met inside an \hbox.
%
% Measured against pdftex. Most of them close the box: the reference says
%
%   ! Missing } inserted.
%   <inserted text>
%                   }
%   <to be read again>
%                      \vfill
%
% and reads the command again, so it does what it came to do in the mode
% outside. `\hbox{A\vfill B}' therefore sets a box holding the A alone,
% then obeys the \vfill in vertical mode, then starts a paragraph for the
% B -- and the source's own } is left over, drawing "Too many }'s".
%
% \vskip and its four fills, \unvbox, \unvcopy and \halign all do this.
%
% Two do not:
%
%   \hrule  is named and ignored, and the horizontal list carries on:
%           `! You can't use `\hrule' here except with leaders.'
%           with `To put a horizontal rule in an hbox or an alignment, /
%           you should use \leaders or \hrulefill (see The TeXbook).'
%           `\hbox{A\hrule B}' sets both letters and no rule.
%   \par    is ignored in silence. There is no paragraph to end inside an
%           \hbox and no shape to clear, so nothing is said at all.
%
% \unhbox is not a vertical command and is unaffected. All of these
% stopped this engine. trip line 377 opens \hbox{\vfill\vsplit 3 0pt}.
\tracingonline=1\showboxbreadth=9999 \showboxdepth=9999
\font\tt=cmr10 \tt
\setbox0=\hbox{A\vfill B}\showbox0
\setbox1=\hbox{A\vskip3pt B}\showbox1
\setbox9=\vbox{}\setbox2=\hbox{A\unvbox9 B}\showbox2
\setbox3=\hbox{A\halign{#\cr X\cr}B}\showbox3
\setbox4=\hbox{A\hrule B}\showbox4
\setbox5=\hbox{A\par B}\showbox5
\setbox8=\hbox{Q}\setbox6=\hbox{A\unhbox8 B}\showbox6
\setbox7=\vbox{A\vskip3pt B}\showbox7      % a paragraph ends instead
\end
